\documentclass[border=4pt]{standalone}
% Shared extension of the installed skill's publication-grade preamble.
% TraceBank figure setup: packages, colour values and plot defaults.
\usepackage{tikz,pgfplots,amsmath,siunitx}
\pgfplotsset{compat=1.18}
\usetikzlibrary{positioning,calc,arrows.meta,fit,backgrounds,shapes.geometric}
\definecolor{cbBlue}{RGB}{0,114,178}
\definecolor{cbOrange}{RGB}{230,159,0}
\definecolor{cbGreen}{RGB}{0,158,115}
\definecolor{cbVermillion}{RGB}{213,94,0}
\definecolor{cbPurple}{RGB}{204,121,167}
\pgfplotsset{every axis/.append style={font=\small,line width=.6pt,
 tick style={black},grid=major,grid style={gray!25,line width=.3pt},
 legend style={font=\footnotesize,draw=gray!50,fill=white},
 label style={font=\small}}}

\usepackage{fontspec}
\usepackage{xeCJK}
\setmainfont{Arial}
\setCJKmainfont{Microsoft YaHei}
\setsansfont{Arial}
\setCJKsansfont{Microsoft YaHei}
\definecolor{ink}{HTML}{17343E}
\definecolor{muted}{HTML}{6C7C80}
\tikzset{
  every node/.style={font=\fontsize{8.3}{12}\selectfont,text=ink},
  box/.style={draw=cbBlue!45,fill=cbBlue!5,rounded corners=2mm,
    minimum height=11mm,text width=31mm,align=center,inner sep=3mm,line width=.6pt},
  safe/.style={box,draw=cbGreen!65!black,fill=cbGreen!7},
  caution/.style={box,draw=cbOrange!80!black,fill=cbOrange!9},
  strong/.style={box,draw=cbVermillion!80!black,fill=cbVermillion!7},
  neutral/.style={box,draw=muted!45,fill=muted!4},
  arrow/.style={-{Stealth[length=2mm,width=1.5mm]},line width=.65pt,draw=ink!75},
  response/.style={arrow,dashed,draw=cbBlue!70},
  note/.style={align=left,font=\fontsize{7.5}{11}\selectfont,text=muted},
  tinylabel/.style={font=\fontsize{7.2}{9}\selectfont,text=muted,fill=white,inner sep=1mm},
  title/.style={anchor=west,font=\fontsize{13}{17}\selectfont\bfseries,text=ink},
  subtitle/.style={anchor=west,font=\fontsize{8}{11}\selectfont,text=muted},
  boundary/.style={draw=cbBlue!45,dashed,rounded corners=3mm,line width=.65pt}
}

% ===== TUNABLES =====
\def\figwidth{183mm}
\def\diagramtitle{F5 · 从需求到回执的演示故事板}
% Based on implemented web flow and live browser acceptance tests.
\begin{document}
\begin{tikzpicture}[x=10mm,y=10mm]
\useasboundingbox (0,.1) rectangle (18.3,11.8);
\node[title] at (0,11.5) {\diagramtitle};
\node[subtitle] at (0,10.9) {已实现的业务交互示意；没有把演示页面包装为用户研究或可用性结果};
\node[safe,text width=45mm,minimum height=29mm] (a) at (3.05,8.5) {\textbf{01\quad 绿色：看懂账单}\\[2mm]“看看本月为什么花多了”\\分类汇总与异常理由\\可展开核对来源交易};
\node[caution,text width=45mm,minimum height=29mm] (b) at (9.15,8.5) {\textbf{02\quad 黄色：确认转账}\\[2mm]“给李娜转 300 元”\\展示收款人、金额与备注\\用户点选明确确认};
\node[safe,text width=45mm,minimum height=29mm] (c) at (15.25,8.5) {\textbf{03\quad 回执：结果有据}\\[2mm]后端完成原子记账\\显示业务 ID 与时间\\失败或未知不伪造成功};
\node[strong,text width=45mm,minimum height=29mm] (d) at (3.05,4.35) {\textbf{04\quad 红色：卡片挂失}\\[2mm]“尾号 1024 的卡丢了”\\明确具体卡片与风险\\不能因检测异常自动挂失};
\node[strong,text width=45mm,minimum height=29mm] (e) at (9.15,4.35) {\textbf{05\quad 独立面板验证}\\[2mm]读取本次操作的模拟 OTP\\验证码不进入聊天上下文\\验证成功后执行卡片操作};
\node[box,text width=45mm,minimum height=29mm] (f) at (15.25,4.35) {\textbf{06\quad 审计：过程可查}\\[2mm]关联计划、授权与回执\\显示哈希链校验结论\\人工接管保留处理记录};
\draw[arrow] (a) -- (b);
\draw[arrow] (b) -- (c);
\draw[arrow] (c.south) -- ++(0,-.55) -- (3.05,6.35) -- (d.north);
\node[tinylabel] at (9.15,6.35) {用户追加需求；不是系统擅自扩展授权};
\draw[arrow] (d) -- (e);
\draw[arrow] (e) -- (f);
\node[neutral,text width=171mm,minimum height=12mm] at (9.15,1.5) {演示边界：全部资金、账户与 OTP 均为模拟；默认离线规则基线。可选模型规划器没有完成实时模型验证时明确披露。\\演示可以增加注入、重复请求和提交后超时，但每次都以实际业务状态和审计记录为依据。};
\end{tikzpicture}
\end{document}
